\documentclass[conference]{IEEEtran}
\IEEEoverridecommandlockouts
\usepackage{cite}
\usepackage{amsmath,amssymb,amsfonts}
\usepackage{algorithmic}
\usepackage{graphicx}
\usepackage{textcomp}
\usepackage{xcolor}
\usepackage{booktabs}
\def\BibTeX{{\rm B\kern-.05em{\sc i\kern-.025em b}\kern-.08em
    T\kern-.1667em\lower.7ex\hbox{E}\kern-.125emX}}

\begin{document}

\title{Transient, SAE-Guided Causal Feature Editing for Factual Correction in Language Models: Architecture, Hypotheses, and Early Empirical Results}

\author{
\IEEEauthorblockN{Chris Walker}
\IEEEauthorblockA{
\textit{Department of Computer Science} \\
\textit{Independent Research Project}\\
chrisiswalker5022@gmail.com}
}

\maketitle

\begin{abstract}
Large language models (LLMs) frequently possess correct factual knowledge internally while producing incorrect outputs, a phenomenon attributable to competition among internal activation directions rather than an absence of the underlying knowledge \cite{elhage2022superposition, meng2022rome}. Existing correction strategies either retrain the model, apply permanent weight edits such as ROME or MEMIT \cite{meng2022rome}, or perform coarse, undirected activation steering. Permanent weight-editing methods have been shown to degrade catastrophically once many edits accumulate \cite{gupta2024forgetting}, while unfiltered activation steering with sparse autoencoder (SAE) features has been shown to steer unreliably unless features are selected by their causal output effect rather than their input activation pattern \cite{arad2025steering}. This paper documents the architecture, hypotheses, and empirical progress of an ongoing project, internally named FeatureScalpel, that uses a pretrained SAE as a diagnostic instrument to identify causally important residual-stream features in GPT-2-small and applies transient, per-inference edits to those features without modifying model weights. Across a 22-fact evaluation subset drawn from the ROME known\_1000 dataset, target-feature ablation succeeded on 4.55\% of suppressed facts, while a corrected strategy that instead ablates the competitor token's driving feature succeeded on 13.64\% of facts (a threefold improvement) with 100\% specificity on held-out control facts. A case study on a persistently unresolved fact revealed that generic grammatical competitor tokens are supported by many overlapping SAE features rather than a single dominant one, a finding we interpret as an empirical instance of the superposition phenomenon predicted by toy models of neural network representation \cite{elhage2022superposition}. A subsequent multi-feature, weighted competitor-reduction strategy partially mitigated this bottleneck. We also report a 21-fold reduction in per-evaluation forward passes (from 122 to 6) achieved through GPU batching while preserving bit-exact equivalence of selected features and safety decisions, yielding up to a 20.67x wall-clock speedup on consumer GPU hardware. We conclude by situating these results against the cited base literature and outlining the specific methodological gaps that remain before a rigorous three-way comparison against weight editing and mean-difference activation steering can be reported.
\end{abstract}

\begin{IEEEkeywords}
sparse autoencoders, mechanistic interpretability, superposition, activation steering, causal feature attribution, knowledge editing, GPT-2
\end{IEEEkeywords}

\section{Introduction}

Large language models are known to sometimes fail to produce a fact they otherwise encode internally: the correct token is present in the model's output distribution but is outranked by a competing, incorrect token. Meng et al. showed that this class of error can be studied causally by tracing the average indirect effect of hidden-state restorations through a corrupted forward pass, and localized factual-association storage to a narrow band of middle-layer multilayer perceptron (MLP) modules acting on the last token of the subject in GPT-2-XL and GPT-J \cite{meng2022rome}. Building on this, Rank-One Model Editing (ROME) directly rewrites the relevant MLP weight matrix to insert a new fact as a closed-form rank-one update \cite{meng2022rome}.

Weight-level editing, however, has a scalability problem. Gupta et al. showed that sequentially applying thousands of ROME or MEMIT edits to the same model produces a two-phase failure mode: a gradual phase in which edit specificity and downstream task performance degrade continuously, followed by an abrupt, catastrophic phase in which the model's editability and general competence collapse, coincident with a sudden spike in the normalized distance between the edited layer's weights and their original values \cite{gupta2024forgetting}. This motivates the central premise of the project reported here: rather than permanently rewriting model weights, we ask whether a factual error can be corrected for a single inference by modifying only the residual-stream activations at inference time, using a sparse autoencoder (SAE) as a diagnostic instrument to identify which directions in activation space are responsible for the error, and discarding the edit after the forward pass completes.

This approach depends on being able to find human- and machine-interpretable, causally meaningful directions in a transformer's activation space. Individual neurons are frequently polysemantic, responding to multiple, semantically unrelated concepts, a phenomenon that Elhage et al. explain with the superposition hypothesis: because real-world features are typically sparse, a network can represent more features than it has neurons by encoding them as an overcomplete set of only approximately orthogonal directions, tolerating the resulting cross-feature interference \cite{elhage2022superposition}. Elhage et al. explicitly propose recovering these superposed directions post hoc as a sparse dictionary-learning problem \cite{elhage2022superposition}, a proposal realized by Cunningham et al. and, independently and concurrently, by Bricken et al., both of whom train sparse autoencoders on a language model's internal activations and show that the resulting dictionary features are substantially more monosemantic and more causally precise than either raw neurons or principal-component directions \cite{cunningham2023sparse, bricken2023towards}. The system described in this paper (internally named FeatureScalpel) uses exactly this class of pretrained SAE as its diagnostic layer.

A further methodological concern, raised specifically in the context of steering rather than diagnosis, is that an SAE feature's activation pattern (what text causes it to fire) is not the same as its causal output effect (what happens to the model's predictions when the feature is amplified or suppressed). Arad, Mueller, and Belinkov show that these two properties are largely disjoint across a feature population and that filtering candidate steering features by an explicit output score, rather than by activation alone, is necessary to obtain reliable steering behavior, closing most of the performance gap between SAE-based steering and supervised representation-editing baselines on the AxBench Concept500 benchmark \cite{arad2025steering}. The causal feature selector at the core of FeatureScalpel is designed around the same principle: features are ranked and selected by their measured causal effect on the target or competitor token's probability under ablation, not by their raw activation magnitude, as elaborated in Section IV.

The remainder of this paper is organized as follows. Section II summarizes the base literature and situates the project relative to it. Section III describes the system architecture as currently implemented in the project repository. Section IV describes the causal feature selection and editing methodology in detail. Section V describes the experimental setup and datasets. Section VI reports empirical results obtained to date. Section VII discusses these results in light of the cited literature. Section VIII states current limitations, and Section IX outlines the remaining work needed to complete the project's stated three-way comparison against weight editing and undirected activation steering.

\section{Related Work}

\subsection{Locating and Editing Factual Knowledge}

Meng et al. introduced causal tracing, a mediation-analysis technique that corrupts a subject's token embeddings with Gaussian noise and then measures how much of a clean prediction is restored by patching individual hidden states back to their clean values, and used it to localize factual recall to middle-layer MLP modules at the last subject token in autoregressive GPT models \cite{meng2022rome}. Their proposed edit, ROME, treats the relevant MLP projection matrix as a linear associative key-value memory and computes a closed-form rank-one update that inserts a new key-value association while adding a KL-divergence penalty to preserve the subject's other properties. On the CounterFact benchmark that the paper introduces, ROME achieves a harmonic-mean score of 89.2 on GPT-2-XL and 91.5 on GPT-J, outperforming fine-tuning and hypernetwork baselines specifically on the joint criterion of generalization and specificity, though the authors note that ROME edits one fact at a time and is directional (editing a forward relation does not automatically edit its inverse) \cite{meng2022rome}.

Gupta, Rao, and Anumanchipalli directly examine the scalability of ROME, MEMIT, and other editing methods by applying up to 2,000 sequential single-fact edits to GPT-2-XL and GPT-J on CounterFact \cite{gupta2024forgetting}. They report that ROME's editing efficacy remains near 100\% until a per-sample inflection point (observed between roughly 100 and 1,000 edits across samples), after which the model suffers catastrophic forgetting of previously edited facts and a corresponding collapse in downstream GLUE task performance; MEMIT delays this catastrophic point substantially (to roughly 1,400 edits on GPT-J in their main sample) at the cost of a lower single-edit efficacy. Both methods show a normalized layer-weight distance that grows gradually and then spikes sharply at the catastrophic point, which the authors interpret as the edited layer's weights becoming incompatible with the rest of the network \cite{gupta2024forgetting}. This result is a central motivation for the transient, non-weight-modifying editing strategy adopted in the present project.

\subsection{Superposition and Sparse Dictionary Learning}

Elhage et al. study small ReLU toy networks trained to reconstruct synthetic sparse input features and show that when the number of features exceeds the number of available hidden dimensions, features are represented non-orthogonally in superposition rather than each being assigned a dedicated dimension, with a sharp, phase-transition-like boundary (as a function of feature sparsity and relative importance) separating the two regimes \cite{elhage2022superposition}. They further show that superposed features organize into specific geometric configurations (antipodal pairs, triangles, tetrahedra, and other polytopes describable as tegum products) and that at least some computation, not merely storage, can be carried out on features while they remain in superposition. The paper explicitly frames the recovery of superposed features as a sparse dictionary-learning problem, directly motivating subsequent SAE-based interpretability work \cite{elhage2022superposition}.

Cunningham et al. train sparse autoencoders, consisting of a single ReLU hidden layer with tied encoder and decoder weights trained with a combined L2-reconstruction and L1-sparsity loss, on the residual stream and MLP activations of Pythia-70M and Pythia-410M \cite{cunningham2023sparse}. They show that the learned dictionary features achieve higher automated interpretability scores than neurons, principal components, or independent-component-analysis directions across all residual-stream layers tested, and that on the indirect-object-identification (IOI) task, patching a small number of dictionary features reaches a target level of output KL divergence with fewer patches and a smaller edit magnitude than patching an equivalent number of principal components \cite{cunningham2023sparse}. Bricken et al. independently obtain analogous results training on a one-layer transformer, additionally reporting that automated interpretability of logit-lens-predicted tokens for learned features reaches 74\% accuracy versus 58\% for neurons, and that pinning a feature to a high activation value during generation reliably produces text matching that feature's inferred concept (for example, forcing generation of Arabic script, DNA sequences, or base64 strings) \cite{bricken2023towards}. Both papers report that reconstructing a layer's activations from the learned dictionary and continuing the forward pass recovers a large majority, but not all, of the log-likelihood the original activations would have provided, indicating that current SAE dictionaries do not yet capture a network's full feature set \cite{cunningham2023sparse, bricken2023towards}.

\subsection{Selecting SAE Features for Causal Intervention}

Arad, Mueller, and Belinkov show that an SAE feature's tendency to activate on a given input token (an input property, measurable via the logit lens without any intervention) and its causal effect on the model's output when amplified (an output property, measurable only via an actual intervention) are largely distinct properties that rarely co-occur strongly in the same feature, and that this distinction, rather than an inherent limitation of SAEs, explains prior reports that SAE-based steering underperforms simpler supervised baselines \cite{arad2025steering}. Filtering candidate steering features by a computed output score before steering raises mean generation-success at $k=20$ from roughly 0.5 to 0.6 in the unfiltered case to 1.1 to 1.4 after filtering, and raises the harmonic-mean AxBench Concept500 score for Gemma-2-9B from 0.293 to 0.546 at layer 20, closing most of the gap to supervised methods such as LoReFT (0.777) \cite{arad2025steering}. This result is directly relevant to the causal feature selector described in Section IV, which likewise ranks candidate features by a measured causal (ablation) effect rather than by activation magnitude, though the present project applies this principle to factual-error correction via ablation and boosting rather than to open-ended concept steering.

\section{System Architecture}

The system, internally named FeatureScalpel within the project repository, is built around GPT-2-small, loaded as a HookedTransformer via the TransformerLens library, together with a pretrained SAE from the SAELens ``gpt2-small-res-jb'' release, which provides a 24,576-feature dictionary attached to the pre-MLP residual stream at each of GPT-2-small's twelve transformer blocks via the hook point \texttt{blocks.\{layer\}.hook\_resid\_pre}. The project's architecture documentation states explicitly that the SAE is used only as a diagnostic instrument for locating and reading out causally relevant directions; the edit itself is applied directly to the residual-stream activation vector, not to the SAE's own weights or to the base model's weights, and is discarded once the forward pass completes.

The processing pipeline proceeds in the following stages, as documented in the repository's architecture description and implemented across the \texttt{src/} package. First, an eligibility scan (\texttt{classify\_fact}) determines whether a given (prompt, target) pair is already correct, suppressed (the correct token is not top-ranked but retains non-trivial probability or a rank within the top 50), or absent from the model's knowledge, since only suppressed facts are candidates for correction. Second, a clean baseline forward pass is cached (the \texttt{CleanContext} data structure), recording the unedited token probabilities, ranks, and the residual-stream activation at the last token. Third, the cached activation is encoded through the SAE to obtain per-feature activations. Fourth, a causal feature selector ranks candidate features, either the features most active on the prompt, or, in the corrected strategy described in Section IV, the features most responsible for the currently predicted competitor token, by their measured effect on the relevant token's probability under full ablation. Fifth, safety filters (\texttt{check\_target\_safe}, \texttt{check\_boost\_safe}, and, for joint edits, \texttt{check\_combination\_safe}) verify that a candidate edit does not regress the target token's rank or probability and does not elevate any previously lower-ranked token above the target (a ``new blocker''). Sixth, the selected features are applied to the residual stream through a forward hook and the remaining transformer layers execute normally on the edited activation, producing the corrected output logits.

The repository implements three deployment surfaces on this pipeline: a minimal demonstration application, a larger ten-tab experimentation and benchmarking interface exposing each intervention strategy (single-feature ablation, compound batch ablation, target boosting, hybrid mute-and-boost, safety-filtered iterative ablation, weighted multi-competitor and multi-feature competitor reduction, and a monosemanticity-analysis tab implementing the max-activating-example and automated-interpretability methodology of Bricken et al. as a diagnostic baseline \cite{bricken2023towards}), and a standalone layer intervention benchmark dashboard that sweeps the intervention across multiple transformer layers and records detailed per-stage timing and forward-pass-count metrics to disk.

\section{Methodology}

\subsection{Residual-Stream Editing via SAE Features}

Following the standard sparse-autoencoder formulation used in the cited base papers \cite{cunningham2023sparse, bricken2023towards}, a feature activation vector is computed as
\begin{equation}
\text{features} = \text{ReLU}(W_{enc} \cdot a + b_{enc})
\end{equation}
from a residual-stream activation $a$, and the SAE's reconstruction of that activation is
\begin{equation}
\text{reconstruction} = W_{dec} \cdot \text{features} + b_{dec}.
\end{equation}
Rather than editing through this full reconstruction, which would inject the SAE's own reconstruction error into the model (since SAE reconstructions are not lossless \cite{cunningham2023sparse, bricken2023towards}), the system computes only a delta in feature space and adds that delta, expressed in the residual-stream basis via the corresponding decoder column, back onto the original, unmodified activation. An ablation of feature $i$ with strength $s \in [0,1]$ is therefore
\begin{equation}
\text{new\_residual} = \text{original\_residual} - s \cdot \text{feature\_activation}[i] \cdot W_{dec}[i],
\end{equation}
and a boost of feature $i$ is
\begin{equation}
\text{new\_residual} = \text{original\_residual} + s \cdot W_{dec}[i].
\end{equation}
Joint edits over a set of mute and boost features are applied within a single hook so that all deltas are computed relative to the same clean baseline.

\subsection{Causal Feature Selection}

Given a prompt and target token, the causal selector does not rank candidate features by their raw activation on the prompt. Instead, for each of the top-$N$ active candidate features it performs a full ablation (strength 1.0) and measures the resulting change in the target (or competitor) token's probability, ranking features by this measured causal delta. An early validation case reported in the project's research notes illustrates why this distinction matters: for the prompt ``The Eiffel Tower is in the city of'' with target `` Paris'', the feature with the highest raw activation on the prompt (feature 5856, activation 34.11) produced a probability change of only $-0.0073$ when ablated, whereas a different feature (feature 11149, activation 30.13, lower than feature 5856) produced a probability change of $-0.0337$ when ablated, roughly 4.6 times larger. This is treated in the project as an internal replication of the general finding that a feature's activation pattern and its causal output effect are separable properties that must be measured independently, echoing the input-score/output-score distinction formalized by Arad et al. for the steering setting \cite{arad2025steering}.

\subsection{Target-Ablation and Competitor-Ablation Strategies}

Two intervention strategies were tested, both operating iteratively for up to five rounds at a fixed ablation strength of 0.3. The first strategy (H1) identifies the feature most causally responsible for the correct target token and ablates it, on the hypothesis that suppressing the target's own supporting feature would, counter-intuitively, allow the target to ``emerge'' by removing whatever mechanism the model was using to hold it down. The second strategy (H2) instead identifies, at each round, whichever token the model currently ranks first (the competitor) and ablates the feature most responsible for driving that specific competitor token, repeating against the new top-ranked competitor if the previous round did not yet elevate the target.

\subsection{Weighted Multi-Competitor and Multi-Feature Reduction}

To address cases in which a single competitor-driving feature is insufficient, two further strategies were implemented. Weighted multi-competitor reduction computes, for every token currently ranked above the target, a weight proportional to that token's baseline probability, maps each competitor token to its single best driving feature, and aggregates ablation strength additively across any features shared by multiple competitors (clamped to a maximum strength), applying all resulting edits jointly. Weighted multi-feature competitor reduction generalizes this further by mapping each competitor to its top-$K$ driving features (rather than only its single best feature) and distributing ablation weight across those $K$ features using one of three weighting schemes (equal weighting, delta-normalized weighting, or softmax weighting with a temperature parameter), again aggregating and applying all edits in one joint pass with the same combination-level safety check used elsewhere in the pipeline.

\subsection{GPU-Batched Causal Evaluation}

Because causal feature selection and safety filtering each require one ablated forward pass per candidate feature, the naive implementation of a single layer-and-prompt evaluation required 122 total model forward passes (1 clean baseline, 30 for competitor-feature ranking, 30 for target-feature ranking, 30 for competitor-feature safety checks, 30 for target-feature safety checks, and 1 for the final intervention). A batched primitive (\texttt{make\_per\_row\_scale\_hook}) repeats the same prompt $B$ times along the batch dimension and scales a distinct candidate feature in each row of the batch during a single forward pass, and on-device rank computation avoids a full-vocabulary argsort. Applying this batching to feature ranking and to safety filtering reduces the same evaluation to 6 total forward passes, a 21-fold reduction, while an internal verification pass reported zero feature-ranking misorderings across six ROME prompts and 100\% agreement (60 of 60 cases) between sequential and batched safety-filter decisions on a canonical case, indicating the optimization is exact rather than approximate.

\section{Experimental Setup}

All experiments use GPT-2-small (117M parameters) with the pretrained ``gpt2-small-res-jb'' SAE release, attached at the pre-MLP residual stream (\texttt{hook\_resid\_pre}) of a specified transformer block. The evaluation fact set is derived from the ROME known\_1000 dataset \cite{meng2022rome}: a 46-fact bank (\texttt{FACT\_BANK}) is used for the eligibility scan and the H1/H2 iterative-ablation experiments described below, of which the subset classified as suppressed (the target token is not the model's top prediction but retains meaningful probability mass or a rank within the top 50) numbers 22 facts and forms the primary evaluation set for Sections VI-A and VI-B. A partially overlapping standardized 22-prompt subset (\texttt{benchmark\_test\_prompts.json}) is used for the layer intervention benchmark reported in Section VI-D. Specificity, following the convention established by Meng et al. for their neighborhood-score metric \cite{meng2022rome}, is measured on a held-out set of already-correct control facts and reports the fraction whose top-1 prediction is unchanged by an intervention calibrated on a different, suppressed fact.

GPU experiments were run on a single NVIDIA GeForce GTX 1660 Ti (6 GB VRAM) under CUDA 12.4 with PyTorch 2.6.0; a discrepancy exists between this empirically used PyTorch build and the version pinned in the repository's \texttt{requirements.txt} (2.13.0), noted here as a reproducibility caveat and discussed further in Section VIII.

\section{Results}

\subsection{Target-Feature Ablation (H1)}

Ablating the feature most causally responsible for the target token itself, iteratively for up to five rounds at strength 0.3, succeeded on 1 of the 22 suppressed facts (4.55\%), with an average of 3.0 rounds used in the single successful case. Manual inspection of this case (``The Eiffel Tower is in the city of'' $\rightarrow$ `` Paris'', successful via a general city-predicting feature that happened to support the wrong competitor `` London'' more strongly than it supported `` Paris'') indicated that the strategy's rare successes were incidental rather than mechanistic: suppressing the target's own supporting feature generally weakened the target further rather than allowing it to emerge, motivating the pivot to competitor-focused ablation.

\subsection{Competitor-Feature Ablation (H2)}

Ablating the feature most responsible for the model's currently top-ranked (incorrect) competitor token, re-selected at each round if the top competitor changed, succeeded on 3 of the 22 suppressed facts (13.64\%), a threefold improvement over target ablation. All three successes required only a single ablation round, and specificity on held-out control facts was 100\% in all three cases, indicating no measured collateral damage to unrelated facts. The three successful cases (Francesco Castellacci $\rightarrow$ Rome, Giulio Romano $\rightarrow$ Rome, Coco Chanel $\rightarrow$ Paris) share the property that the model's top-ranked incorrect token was a specific, semantically contentful competing fact (for example, a wrong city name) rather than a generic function word.

\subsection{The Grammatical-Competitor Bottleneck}

The remaining 19 of 22 cases failed under competitor ablation, and inspection showed that in the majority of these the model's top-ranked competitor was a generic grammatical or syntactic token (for example, `` the'', `` a'', or `` an'') rather than a specific wrong fact. A detailed case study on the prompt ``The location of Massachusetts Institute of Technology is in'' (target `` Cambridge'', clean rank 8, clean probability 1.15\%) found that its seven above-target competitor tokens (`` the'' at 17.95\%, `` a'' at 10.53\%, `` question'' at 7.51\%, `` jeopardy'' at 2.71\%, `` an'' at 2.63\%, `` danger'' at 1.76\%, `` doubt'' at 1.21\%) reduced, under a max-of-shared-strength feature attribution, to only two unique driving features (feature 313, driving `` the'', `` a'', and `` an''; feature 21169, driving `` question'', `` jeopardy'', `` danger'', and `` doubt''), and that muting these two features left the target's rank unchanged at 8. We interpret this as a case in which a single SAE feature cannot capture a token's full causal support because that support is genuinely distributed across many more directions in activation space than the dictionary size and ablation strategy used here can individually address, a pattern consistent with the general expectation from superposition theory that frequently occurring, low-information tokens are supported by broad, densely shared representations across many features rather than by one dedicated feature \cite{elhage2022superposition}. This finding was formalized in the project's research log as the Distributed Support Hypothesis.

\subsection{Mitigating Distributed Support: Multi-Feature Reduction}

Testing whether the target token's rank on the same MIT-to-Cambridge case could be improved by distributing ablation weight across each competitor's top-$K$ driving features rather than only its single best feature, using $K=3$ with softmax or equal weighting recovered the target's rank from 8 to 6, whereas summation-based aggregation of shared single-best features alone (rather than top-$K$ distribution) raised the target probability marginally (from 1.29\% to 1.39\%) without improving rank, and increasing $K$ to 5 was found to dilute the edit and reduce its effectiveness. This indicates that the grammatical-competitor bottleneck identified in Section VI-C is only partially addressed by widening the feature set considered per competitor, and that rank 1 (a full correction) was not achieved on this case within the strategies tested.

\subsection{Layer Intervention Benchmark and Compute Efficiency}

Sweeping the hybrid mute-and-boost intervention (competitor muting combined with target boosting, applied jointly with safety filtering) across transformer blocks $\{2, 5, 8, 10\}$ on the MIT-to-Cambridge case found layer 8 to be the most effective of the four layers tested, improving the target's rank from 8 to 3 (a gain of 5 rank positions) and its probability by 2.37 percentage points, compared to smaller gains at layers 2 (rank 8 to 7, +0.07 points), 5 (rank 8 to 5, +0.79 points), and 10 (rank 8 to 5, +0.82 points), as summarized in Table \ref{tab:layer}. Note that in this specific recorded run the model's top-1 prediction did not flip to the target token even though its rank and probability improved, illustrating that ``success'' as measured here (an improvement in rank or probability) is a weaker criterion than full correction of the model's top-1 output.

\begin{table}[htbp]
\caption{Layer Sweep on the MIT $\rightarrow$ Cambridge Case (Hybrid Mute-and-Boost, Clean Rank 8)}
\label{tab:layer}
\centering
\begin{tabular}{@{}cccl@{}}
\toprule
\textbf{Layer} & \textbf{Rank (clean $\to$ final)} & \textbf{Prob. gain (pts)} & \textbf{Notes} \\
\midrule
2  & 8 $\to$ 7 & +0.07 & Weakest effect \\
5  & 8 $\to$ 5 & +0.79 & Moderate \\
8  & 8 $\to$ 3 & +2.37 & Best of 4 layers tested \\
10 & 8 $\to$ 5 & +0.82 & Comparable to layer 5 \\
\bottomrule
\end{tabular}
\end{table}

The GPU-batched forward-pass optimization described in Section IV-E reduced per-layer-per-prompt evaluation from 122 to 6 total forward passes. On a fixed three-prompt, four-layer evaluation (prompts concerning Dublin, Rome, and Paris facts) run on the GTX 1660 Ti GPU versus the same evaluation on CPU, the project's speed-analysis script recorded a GPU speedup ranging from 3.83x to 20.67x, corresponding to approximately 1.63 seconds per prompt on GPU versus 33.62 seconds per prompt on CPU in the best case, with a peak GPU memory utilization of 1.45 GB of the 6.00 GB available.

\section{Discussion}

The overall pattern of results, low but non-zero success under single-feature ablation, a threefold improvement from correctly targeting the competitor rather than the target, and a hard bottleneck on tokens whose support is distributed across many features, is broadly consistent with the theoretical picture established by the cited base literature rather than contradicting it. Elhage et al. predict that a network with more features than available dimensions will represent common, low-sparsity features (which, for a language model, plausibly includes frequent function words) using shared, overlapping directions rather than one dedicated direction each \cite{elhage2022superposition}; the two-feature collapse of seven distinct competitor tokens observed in Section VI-C is a direct, if informal, empirical instance of this prediction at the level of a trained SAE dictionary rather than a toy model. Bricken et al. and Cunningham et al. both report that current SAE dictionaries recover a substantial majority, but not all, of a layer's causally relevant information, leaving a residual that is not captured by any single learned feature \cite{cunningham2023sparse, bricken2023towards}; the residual, unascribed causal support behind the grammatical-competitor bottleneck is plausibly related to this same incompleteness, rather than being a distinct algorithmic failure of the causal selector.

The project's causal, ablation-based feature ranking is methodologically aligned with the central argument of Arad et al., that a feature's causal output effect must be measured directly rather than inferred from its activation pattern \cite{arad2025steering}. However, the present system applies this principle only per feature (or, in the multi-feature reduction strategies, to a small, greedily-selected top-$K$ set of features per competitor), whereas the grammatical-competitor case study suggests that some tokens require intervening on a genuinely higher-dimensional subspace of features simultaneously. This gap between single- or few-feature causal ranking and subspace-level intervention is, to date, unresolved in the project and is the most direct link between the empirical bottleneck reported here and the feature-selection literature this project builds on.

With respect to the project's original comparative aim, namely a controlled, apples-to-apples comparison among permanent weight editing (ROME/MEMIT), undirected mean-difference activation steering, and SAE-guided transient activation editing, the results reported here establish only the third arm of that comparison. The motivation drawn from Gupta et al., that weight editing degrades unpredictably at scale \cite{gupta2024forgetting}, remains a stated rationale for the transient approach rather than a result independently reproduced within this project, since no ROME or MEMIT edits have yet been run on the same fact set within the repository as of this report.

\section{Limitations}

Several limitations constrain the generality of the results reported here. First, the evaluation set is small: 22 suppressed facts (drawn from a 46-fact bank) for the H1/H2 iterative-ablation experiments, which limits the statistical precision of the reported success rates. Second, all experiments use a single, comparatively small base model, GPT-2-small (117M parameters), with a single pretrained SAE release; whether the same feature-selection dynamics and grammatical-competitor bottleneck hold in larger or instruction-tuned models is untested within this project. Third, the causal editing pipeline (\texttt{src/editing.py}, \texttt{src/hooks.py}) is currently validated only through ad hoc verification scripts and manual interface inspection rather than an automated regression test suite; automated tests in the repository cover only the interpretability and diagnostic tooling (feature-similarity computation, max-activating-example extraction, and input-validation boundary cases), not the correctness of the ablation, boosting, or safety-filtering logic itself. Fourth, a version discrepancy exists between the PyTorch build pinned in the repository's dependency file (2.13.0) and the CUDA-enabled build actually used for the GPU benchmarking results reported in Section VI-E (2.6.0 with CUDA 12.4), which should be resolved before the reported speedup figures are treated as fully reproducible from a clean checkout. Fifth, the raw run artifacts underlying the 4.55\% and 13.64\% headline success rates (Sections VI-A and VI-B) are recorded in narrative research-journal entries with full per-fact result tables, but the standalone experiment notebooks intended to reproduce them programmatically (\texttt{research/notebook/exp\_001.md}, \texttt{exp\_003.md}) are marked in the repository as pending reconstruction. Finally, ``success'' throughout Sections VI-A, VI-B, and VI-E is defined as an improvement in the target token's rank or probability, which is a weaker criterion than the target becoming the model's top-1 prediction; several reported cases show measurable rank or probability improvement without a full top-1 correction.

\section{Future Work}

The most immediate remaining work, and the comparison the project's own architecture documentation identifies as its core novelty claim, is a controlled evaluation of permanent weight editing (ROME and MEMIT \cite{meng2022rome}), undirected mean-difference activation steering (the baseline strategy used for comparison by Arad et al. and by the AxBench benchmark referenced therein \cite{arad2025steering}), and the SAE-guided transient editing strategy described in this paper, on the same base model, the same fact set, and the same evaluation metrics. Second, the grammatical-competitor bottleneck identified in Section VI-C motivates moving from top-$K$ per-competitor feature selection toward genuinely subspace-based interventions that can represent and edit a distributed, many-feature representation as a single object rather than as an aggregate of individually-selected features. Third, scaling the evaluation to larger and instruction-tuned models, and to a larger and more diverse fact set than the current 22 to 46 ROME-derived facts, is necessary before the reported success rates can be treated as representative rather than as a pilot result. Fourth, given the caution raised by Gupta et al. that weight-editing degradation is only observable at scale \cite{gupta2024forgetting}, an equivalent scaling study for the transient editing approach, repeatedly applying and discarding edits across a large number of sequential facts, would test whether the project's transient design in fact avoids the gradual and catastrophic forgetting dynamics reported for weight editing, a claim that is currently argued from first principles rather than measured directly.

\section{Conclusion}

This paper reported the current architecture and empirical progress of a project investigating whether sparse-autoencoder-identified, causally-ranked residual-stream features can be used to transiently correct suppressed factual knowledge in GPT-2-small without modifying model weights. Correctly targeting the SAE feature driving a model's incorrect competitor prediction, rather than the feature driving the correct target, tripled correction success (from 4.55\% to 13.64\%) on a 22-fact evaluation subset drawn from ROME's known\_1000 dataset, with perfect measured specificity on control facts, while revealing that generic grammatical competitors resist single-feature correction because their causal support is distributed across multiple, overlapping SAE features, a pattern consistent with the superposition phenomenon documented in the toy-model and sparse-autoencoder literature this project builds on \cite{cunningham2023sparse, bricken2023towards, elhage2022superposition}. A GPU-batched evaluation primitive reduced the computational cost of causal feature ranking and safety filtering by a factor of 21 without changing any selection or safety outcome, making a substantially larger comparative study, including the weight-editing and undirected-steering baselines the project was designed to compare against, computationally tractable as the next phase of this work.

\begin{thebibliography}{00}
\bibitem{cunningham2023sparse} H. Cunningham, A. Ewart, L. Riggs, R. Huben, and L. Sharkey, ``Sparse Autoencoders Find Highly Interpretable Features in Language Models,'' arXiv:2309.08600, 2023.
\bibitem{bricken2023towards} T. Bricken, A. Templeton, J. Batson, B. Chen, A. Jermyn, et al., ``Towards Monosemanticity: Decomposing Language Models With Dictionary Learning,'' Transformer Circuits Thread, Anthropic, Oct. 2023.
\bibitem{elhage2022superposition} N. Elhage, T. Hume, C. Olsson, N. Schiefer, T. Henighan, S. Kravec, et al., ``Toy Models of Superposition,'' Transformer Circuits Thread, Anthropic, Sept. 2022.
\bibitem{meng2022rome} K. Meng, D. Bau, A. Andonian, and Y. Belinkov, ``Locating and Editing Factual Associations in GPT,'' in \textit{Proc. 36th Conf. Neural Information Processing Systems (NeurIPS)}, 2022, vol. 35, pp. 17359--17372.
\bibitem{gupta2024forgetting} A. Gupta, A. Rao, and G. Anumanchipalli, ``Model Editing at Scale leads to Gradual and Catastrophic Forgetting,'' arXiv:2401.07453, 2024.
\bibitem{arad2025steering} D. Arad, A. Mueller, and Y. Belinkov, ``SAEs Are Good for Steering, If You Select the Right Features,'' arXiv:2505.20063, 2025.
\bibitem{wu2025axbench} Z. Wu, A. Geiger, A. Wang, J. Huang, C. D. Manning, and C. Potts, ``AxBench: Steering LLMs? Even Simple Baselines Outperform Sparse Autoencoders,'' in \textit{Proc. Int. Conf. Machine Learning (ICML)}, 2025. (Cited as reported and evaluated against in \cite{arad2025steering}.)
\end{thebibliography}

\end{document}
